\subsection{A four-valued \texorpdfstring{$\delta$}{delta}-filter on the fixed-point algebra}
\label{sec:fitting-delta-filter}

Section~\ref{sec:quotient-delta-filter} placed a four-valued \(\delta\)-filter on the abstract quotient algebra, where the filter is an arbitrary family of quotient extensions. The closure construction also supports the corresponding question one level down, on the selected admissible domain itself, where the filter is the positivity predicate \(P\) rather than an auxiliary family. This is the level at which Fitting's extensional \(\delta\)-filter lives, and the level at which maximality can interact with Godlikeness. Stating it here also keeps positivity-filter closure separate from closure of the ambient property domain, the distinction forced by the entailment result above.

The order is the FDE truth order \(\leq_t\) of Section~\ref{sec:quotient-delta-filter}, not the information order \(\leq_k\) used by the profile hull. Lean verifies that the pointwise FDE operations introduced above are meet and join for \(\leq_t\), while FDE negation reverses this order. Thus \(Sat_w\) is a closure operator in the information order, whereas the filter structure on its fixed points is truth-ordered.

I call the selected domain an admissible FDE algebra when it contains \(\top_F,\bot_F\) and is closed under FDE conjunction and disjunction; negation closure is already built into the admissible semantics. A profile \(\delta\)-filter then combines this algebra condition with profile saturation and, at every world, the laws
\[
  +P(\top_F),
  \qquad
  \neg +P(\bot_F),
\]
\[
  +P(X)\land X\leq_tY\Longrightarrow +P(Y),
\]
and
\[
  +P(X)\land +P(Y)\Longrightarrow +P(X\wedge_FY),
\]
for admissible \(X,Y\). The preceding fixed-point theorem then places every carrier extension among the fixed points of \(Sat_w\).

This definition has the intended classical recovery. Let \(X\leq_+Y\) abbreviate \(X^{+}\subseteq Y^{+}\), the ordinary subset order used by the classical \(\delta\)-filter
\cite{BenzmuellerFuenmayor2019}. Truth-order inclusion always entails \(\leq_+\). If admissible extensions are classically coherent, positive inclusion conversely entails truth-order inclusion because negative membership is complementary to positive membership. Lean therefore proves equivalence between the FDE truth-order filter and the source-style positive-set filter on the coherent fragment.

The bilateral extension adds the polarity equation
\[
  -P(X)\Longleftrightarrow +P(\neg_FX).
\]
Lean derives both admissible A1 directions and shows that negative support forms the dual ideal: it contains \(\bot_F\), excludes \(\top_F\), is downward closed in \(\leq_t\), and is closed under FDE disjunction. This equation still permits both positivity gluts and positivity gaps; it is not a bivalence condition.

Ultrafilter maximality remains separate:
\[
  \forall X\in Adm,
  \qquad
  +P(X)\lor +P(\neg_FX).
\]
As this is precisely complement decision, the earlier dependency result applies. Lean proves
\[
\boxed{
  \text{profile }\delta\text{-ultrafilter}
  +CONS_{G,F}^{G,adm}
  \Longrightarrow
  COMP_P^{G,adm},
}
\]
and obtains the same conclusion with \(A1_R\) in place of relevant exemplification consistency. The proof consumes only the maximality clause; the filter and saturation conditions are carried by the definition but not used. Hence maximality returns to the classification route on its own, and it is not part of the genuinely \(COMP\)-independent filter core.

A finite witness separates these notions. Its admissible algebra is the 12-element product
\[
  \{\True,\False,\Neither\}
  \times
  \{\True,\False,\Neither,\Both\}
\]
over two entities. It is closed under the FDE operations and contains both gap and glut information. Positive support is the truth-order filter generated by \((\True,\Neither)\). This generator distinguishes the two entities, so the induced positive-profile quotient is discrete and all admissible extensions are fixed points. The model has a positive Godlike witness and satisfies relevant exemplification consistency. Nevertheless, for the split extension \(A=(\True,\False)\),
\[
  +A(a),
  \qquad
  \neg +P(A),
  \qquad
  \neg -P(A).
\]
Therefore the profile \(\delta\)-filter plus relevant consistency does not imply \(COMP_P^{G,adm}\), while the corresponding ultrafilter does so generally. This strict separation is the main reason to keep filter closure, profile closure, and maximality as three distinct interfaces.
